% Source: https://artofproblemsolving.com/community/c3h3162282_2016_lmt_spring_guts_round_p2536__lexington_mathematical_tournament
% Stable identity: aops-topic-3162282-q30
Write $\phi=\dfrac{1+\sqrt5}{2}$. It suffices to show that at $ k=\phi $ there are fewer than eight roots, while every $\phi<k\le2 $ has eight. On each step, label the branch $ L $ if the input is at most $\dfrac{1}{2}$ and $ R $ if it is at least $\dfrac{1}{2}$. Composing the affine maps $ L(u)=ku $ and $ R(u)=k-ku $ and solving the fixed-point equation gives, in itinerary order $ LLL,LLR,LRL,LRR,RLL,RLR,RRL,RRR $:
\[0,\quad\dfrac{k}{k^3+1},\quad\dfrac{k^2}{k^3+1},\quad\dfrac{k}{k^2+k+1},\quad\dfrac{k^3}{k^3+1},\quad\dfrac{k(k+1)}{k^2+k+1},\quad\dfrac{k^2}{k^2+k+1},\quad\dfrac{k}{k+1}.\]
For $\phi<k\le2 $, substitution into their three branch inequalities verifies every listed itinerary; the limiting condition is $ k^2-k-1>0 $. The eight values are distinct, so all eight are roots. There can be no others, as each of the eight affine itineraries has at most one fixed point. At $ k=\phi $, the second and fourth values coincide, the third and seventh equal $\dfrac{1}{2}$, and the fifth and sixth equal $\dfrac{\phi}{2}$. Only five distinct roots remain. Thus the desired maximum is $\boxed{\phi}$.
